%=============================================================================!
% 3.6  ENERGY DIAGRAMS                                                        !
%=============================================================================!
\section{Energy diagrams}
\label{sec:en-diagrams}

The conservation of energy turns a graph of the potential energy into a
map of every possible motion. This section reads such a map.

\subsection*{Allowed regions and turning points}

For a body moving along a line with potential energy $U(x)$, or a bead on
a wire with $U = mgh(x)$, the kinetic energy at each point is what
remains of the total, $K = E - U(x)$. A kinetic energy $\frac12 mv^2$
cannot be negative, so the body can only be where
\begin{equation*}
   U(x) \le E .
\end{equation*}
These points form the allowed region of the motion. Where $U$ rises to
meet $E$, the kinetic energy falls to zero, the body stops for an
instant, and the force, which points down the slope of $U$ by
\cref{eq:en-force-1d} (for the bead, the part of the weight along the
wire, which points downhill), sends it back the way it came. Such a point, where
$U(x) = E$ with a non-zero slope of $U$, is a turning point. If $E$ instead
meets $U$ at an equilibrium, the force also vanishes there: a body placed
there at rest stays there, and approaching an unstable equilibrium can
take indefinitely long. Away from these special points, the speed follows
from $\frac12 mv^2 = E - U(x)$,
\begin{equation*}
   v = \sqrt{\frac{2\,\bigl(E - U(x)\bigr)}{m}} .
\end{equation*}

A graph of $U$ with a horizontal line at the height $E$ is called an
energy diagram. For the bead, $U/mg = h$ is the shape of the wire itself,
so its energy diagram is a picture of the track, and $E/mg = h_E$ is the
height at which the bead would be at rest. \Cref{fig:en-track}(a) draws
three energies. Their turning points, the solutions of $h(x) = h_E$, were
found by computer; substituting them into \cref{eq:en-track} checks them,
and two are exact, $h(-1) = 0.4$ and $h(2) = 1.3$.
\begin{itemize}
\item At $h_E = \SI{0.4}{\metre}$ the bead is confined to the deep valley,
   between the turning points $x = \SI{-1.831}{\metre}$ and
   \SI{-1.000}{\metre}. It slides back and forth between them for ever.
\item At $h_E = \SI{0.8}{\metre}$ there are two separate allowed regions,
   one in each valley: from \num{-2.042} to \SI{-0.477}{\metre}, and from
   \num{0.795} to \SI{1.724}{\metre}. With the same energy the bead may
   be in either, but it cannot cross the hill between them, which rises to
   \SI{1.009}{\metre}. Which valley it occupies depends only on where it
   started.
\item At $h_E = \SI{1.3}{\metre}$ the bead passes over the hill and turns
   only at the steep ends of the track, at \num{-2.206} and
   \SI{2.000}{\metre}.
\end{itemize}
\Cref{fig:en-track}(b) gives the speed along the track at the highest of
these energies. It falls to zero at the turning points, is greatest at the
bottom of the deep valley, \SI{4.680}{\metre\per\second}, and dips to
\SI{2.387}{\metre\per\second} as the bead crosses the hill.

\begin{figure}[tb]
\centering
\includegraphics{ch03_energy/track}
\caption{The track as an energy diagram. (a) The potential energy of the
bead divided by $mg$, which is the height $h(x)$ of the wire, with three
total energies $E/mg$. Each energy line is
solid where the bead may be ($U \le E$) and dotted where it may not; dots
mark the turning points. Filled circles mark the stable equilibria at the
valley bottoms, and the open circle the unstable one at the top of the
hill. (b) The speed along the track at $E/mg = \SI{1.3}{\metre}$; the
dotted lines are the turning points.}
\label{fig:en-track}
\end{figure}

\subsection*{Equilibria}

A body placed at rest stays at rest only where the net force on it is
zero. By
\cref{eq:en-force-1d} these are the points where $\dd U/\dd x = 0$, where
the graph of $U$ is flat: the bottoms of valleys, the tops of hills, and
any level shelf. Such a point $x_0$ is called an equilibrium. For the
bead, the push of the wire and the weight balance where the wire is
horizontal, $\dd h/\dd x = 0$, since only there is the push, which is
perpendicular to the wire, vertical, so that it can cancel the weight. Differentiating \cref{eq:en-track} with the
power rule,
\begin{equation*}
   \frac{\dd h}{\dd x} = 0.6\,x^3 - 1.2\,x + 0.15 ,
\end{equation*}
and this cubic is zero at three points, which a computer finds to be $x =
\SI{-1.473}{\metre}$ (the deep valley, $h = \SI{0.183}{\metre}$), $x =
\SI{0.126}{\metre}$ (the hill, $h = \SI{1.009}{\metre}$) and $x =
\SI{1.347}{\metre}$ (the shallow valley, $h = \SI{0.607}{\metre}$).

The equilibria differ in what happens after a small nudge. At the bottom
of a valley, the slope of $U$ on either side points back down into the
valley, so the force on a displaced body points back towards the
equilibrium, and the body swings about it: the equilibrium is stable. At
the top of a hill, the slope on either side leads away, and a body nudged
off it runs away down one side: the equilibrium is unstable.

The second derivative of $U$ tells the two apart without a picture. Near
an equilibrium $x_0$, the Taylor series of \cref{sec:mo-taylor}, applied
to the force $F$ about $x_0$ with $x - x_0$ in the role of the interval
$\tau$, gives
\begin{align*}
   F(x) &= F(x_0) + F'(x_0)\,(x - x_0) + \order{(x - x_0)^2}\\
   &= -U''(x_0)\,(x - x_0) + \order{(x - x_0)^2}
      && \text{since $F(x_0) = 0$ and $F' = -U''$,}
\end{align*}
where $U''$ is the second derivative of $U$ with respect to $x$. If
$U''(x_0) > 0$, the graph of $U$ curves upwards, as at a valley bottom,
and close to $x_0$ the force is that of a spring of stiffness $k =
U''(x_0)$, $F \approx -k\,(x - x_0)$: it always points back, the
equilibrium is stable. By \cref{sec:ne-spring}, with the displacement
$x - x_0$ in place of $x$, the body then oscillates about $x_0$,
approximately, with angular frequency $\sqrt{U''(x_0)/m}$,
and Chapter~4 examines how good the approximation is. If $U''(x_0) <
0$, the force points away on both sides, and the equilibrium is unstable.
If $U''(x_0) = 0$, the next terms of the series decide. For the track,
$\dd^2 h/\dd x^2 = 1.8\,x^2 - 1.2$, which is \SI{+2.706}{\per\metre} and
\SI{+2.066}{\per\metre} at the two valley bottoms and
\SI{-1.171}{\per\metre} at the top of the hill, as the picture says.

\notebook{03}{move $E$ on the energy diagram and watch the allowed regions
change}

\begin{selfcheck}
At $E/mg = \SI{0.8}{\metre}$ a bead starts at rest at $x =
\SI{0.795}{\metre}$. Where does it next turn back, where is it fastest,
and how fast is it moving there?
\end{selfcheck}
